%% REVTeX 4-2 — Physical Review E
%% Compile: pdflatex → bibtex → pdflatex × 2
%% Figures must be in the same directory as this .tex file
\documentclass[aps,pre,twocolumn,showpacs,superscriptaddress,longbibliography]{revtex4-2}

\usepackage{amsmath,amssymb,amsfonts}
\usepackage{graphicx}
\usepackage{dcolumn}
\usepackage{bm}
\usepackage{hyperref}
\usepackage{xcolor}
\usepackage{booktabs}
\usepackage{multirow}

%% ── Custom commands ─────────────────────────────────────────
\newcommand{\lp}{\lambda_{+}}
\newcommand{\lmax}{\lambda_{\max}}
\newcommand{\rGOE}{\langle r \rangle}
\newcommand{\kN}{k/N}
\newcommand{\taupop}{\tau_{\mathrm{pop}}}
\newcommand{\tauISI}{\tau_{\mathrm{ISI}}}
\newcommand{\tauabs}{\tau_{\mathrm{abs}}}
\newcommand{\hneu}{h_{\mathrm{neural}}}
\newcommand{\rhoinfo}{\rho_{\mathrm{info}}}
\newcommand{\fUV}{f_{\mathrm{UV}}}

\begin{document}

%% ════════════════════════════════════════════════════════════
%%  TITLE & AUTHORS
%% ════════════════════════════════════════════════════════════
\title{Power-law spectral structure of synaptic connectivity matrices
       is conserved across cortical areas, layers, and species}

\author{Gesivaldo Santos}
\email{gesivaldo@uesb.edu.br}
\affiliation{Department of Biological Science,
             State University of Southwest Bahia (UESB),
             Jequié, Bahia, Brazil}

\date{\today}

%% ════════════════════════════════════════════════════════════
%%  ABSTRACT
%% ════════════════════════════════════════════════════════════
\begin{abstract}
The cortical microcircuit operates under strict metabolic constraints
that limit simultaneous neuronal firing to approximately 15\% of the
local population. We show that this constraint, formalized within a
thermodynamic framework for neural ensemble dynamics, predicts a
discrete set of collective synaptic modes---attractors---whose spectral
structure follows a universal power law. Applying random matrix theory
to two independent connectivity datasets---the Blue Brain Project
neocortical simulation (BBP, $n=8{,}173$ neurons, rat somatosensory
cortex L2/3) and the MICrONS electron microscopy reconstruction
($n=50{,}948$ neurons, mouse visual cortex, measured synaptic volumes
in nm$^3$)---we find that the ranked eigenvalues of the symmetrized
synaptic weight matrix follow $\lambda_k \sim k^\alpha$ with
$\alpha = {-}0.450 \pm 0.074$ (BBP, 95\% CI $[-0.539, -0.260]$)
and $\alpha = {-}0.422 \pm 0.062$ (MICrONS, 95\% CI $[-0.502, -0.267]$),
statistically indistinguishable across species, cortical areas, and
data modalities (bootstrap $n = 1{,}000$; overlapping CIs).
This exponent is destroyed by degree-preserving rewiring
(null model $\alpha = {-}0.141 \pm 0.009$, $p < 0.001$,
$n = 500$ permutations), and is the preferred model over exponential
and harmonic alternatives by formal AIC selection ($\Delta\mathrm{AIC} > 23$
in both datasets). Both datasets exhibit $\sim$20 eigenvalues above the
Marchenko-Pastur bulk, with level-spacing statistics consistent with the
Gaussian Orthogonal Ensemble ($\rGOE = 0.535$ and $0.654$).
A laminar analysis across cortical layers reveals a monotonic depth
gradient in $\alpha$ (L2/3: $-0.451$ to L5/L6: $-0.531$,
$R^2 = 0.91$), confirmed by a layer-matched cross-dataset comparison
(BBP L4: $\alpha = -0.525$; MICrONS L4: $\alpha = -0.557$,
$\Delta\alpha = 0.032$, within bootstrap uncertainty).
We derive from first principles that the maximum number of
non-overlapping active ensembles before universal refractoriness is
$n_{\max} = \lfloor N/k \rfloor = 6$, providing a thermodynamic
derivation of the six canonical EEG frequency bands
($f_n = 1/(n\taupop)$, $n = 1,\ldots,6$). Two distinct spectral
cutoffs emerge: an ensemble cutoff (Debye analog,
$\fUV^{\mathrm{ens}} = 114\,\mathrm{Hz}$ at rest) and an
invariant neuronal cutoff (Planck analog,
$\fUV^{\mathrm{neural}} = 587\,\mathrm{Hz}$), consistent with
empirical reports of very high-frequency oscillations in human
somatosensory cortex (407--926\,Hz) and thalamus (1235--2632\,Hz).
Phase-space reconstruction via Takens delay embedding reveals a
quasi-1D attractor ($d_f \approx 1.1$) with negative maximal Lyapunov
exponent ($\lambda_L = {-}0.224$ in simulation, $-0.082$ in EM data),
with the EM value approaching the critical point ($\lambda_L \to 0$),
consistent with the cortical criticality hypothesis.
\end{abstract}

\pacs{87.19.lj, 87.19.lv, 05.45.Tp, 87.10.Mn}

\keywords{random matrix theory, synaptic connectivity, power-law spectrum,
Marchenko-Pastur, Wigner-Dyson, electron microscopy, Blue Brain Project,
MICrONS, cortical attractors, neural thermodynamics, Takens embedding,
UV catastrophe, EEG frequency bands}

\maketitle

%% ════════════════════════════════════════════════════════════
%%  I. INTRODUCTION
%% ════════════════════════════════════════════════════════════
\section{\label{sec:intro}Introduction}

The neocortical microcircuit---a canonical column of approximately
$N \approx 100$ neurons~\cite{Mountcastle1997}---constitutes the
fundamental computational unit of the mammalian cerebral cortex.
A central empirical constraint on its operation is the sparse activation
regime: at any moment, only a fraction $\kN \approx 0.15$ of the
local population fires simultaneously, a constraint arising from
the strict metabolic budget of cortical tissue~\cite{Attwell2001}.

Random matrix theory (RMT) provides a natural framework for characterizing
the spectral properties of complex connectivity matrices~\cite{Mehta2004}.
The bulk of the eigenvalue spectrum of a random synaptic weight matrix is
described by the Marchenko-Pastur (MP) distribution. Eigenvalues exceeding
the MP upper edge $\lp$ indicate structured, low-rank components---the
spectral signature of learned or evolved synaptic
attractors~\cite{Mastrogiuseppe2018}.

Here we combine RMT analysis with two large-scale connectivity datasets:
the Blue Brain Project (BBP) neocortical microcircuit model~\cite{Markram2015},
and the MICrONS dataset~\cite{MICrONS2021}, a volumetric electron microscopy
(EM) reconstruction of mouse visual cortex providing direct measurements
of synaptic volumes in nm$^3$---a physical proxy for synaptic strength.

We find that the outlier eigenvalues in both datasets follow a universal
power law $\lambda_k \sim k^\alpha$ with $\alpha \approx -0.43$,
conserved across species, cortical areas, and data modalities. Formal
AIC model selection confirms the power law as the preferred model
($\Delta\mathrm{AIC} > 23$). We further derive the discrete EEG
frequency bands and two UV spectral cutoffs from first principles
using metabolic constraints alone, and show that the phase-space
geometry of the spectral attractor is consistent with cortical
criticality.

%% ════════════════════════════════════════════════════════════
%%  II. METHODS
%% ════════════════════════════════════════════════════════════
\section{\label{sec:methods}Data and Methods}

\subsection{Datasets}

\paragraph{Blue Brain Project (BBP).}
We use the publicly available neocortical microcircuit model
of juvenile rat somatosensory cortex from the Open Brain Institute
AWS S3 repository (\texttt{s3://openbluebrain/}).
The L2/3 matrix: $n = 8{,}173$ neurons, $974{,}129$ synapses.
We also analyzed layers L4 ($n=4{,}870$), L5 ($n=5{,}824$),
and L6 ($n=11{,}032$).

\paragraph{MICrONS V1.}
The MICrONS cortical mm$^3$ dataset~\cite{MICrONS2021} provides
EM reconstruction of mouse primary visual cortex (V1):
$n = 50{,}948$ neurons, $11{,}554{,}213$ synapses, with measured
synaptic volumes in nm$^3$ (field \texttt{source\_\_size}).
Log-transformed weights $w_{ij} = \ln(V_{ij}^{\mathrm{syn}})$
are used for spectral analysis, consistent with the empirical
log-normal distribution of synaptic sizes~\cite{Bhatt2009}.
A layer-matched L4 subsample was extracted via cortical depth
coordinates ($n_{L4} = 10{,}188$ neurons,
$1{,}664{,}367$ L4$\to$L4 synapses).

\paragraph{Electrophysiology.}
Nine patch-clamp recordings from rat SSCx neurons (IDRest protocol)
from the BBP electrophysiology repository were analyzed to validate
$\taupop$.

\subsection{Spectral analysis}

The symmetrized weight matrix $J_s = (J + J^T)/2$ was analyzed.
The MP upper edge is $\lp = 4\sigma^2$ (square matrix case).
Top $k=30$ eigenvalues were computed via sparse Lanczos
(\texttt{scipy.sparse.linalg.eigsh}, ARPACK).
The power-law exponent $\alpha$ was fit by OLS in log-log space;
95\% confidence intervals by non-parametric bootstrap ($n=1{,}000$).

\subsection{Model selection}

Four models were compared via AIC ($= 2k - 2\ln\hat{L}$):
power law $\lambda_k \sim k^\alpha$ (2 params),
Boltzmann $\sim e^{Bk}$ (2 params),
harmonic $\sim \lambda_1/k$ (1 param),
and log-normal $\sim e^{a(\ln k)^2}$ (2 params).

\subsection{Null models}

Three levels: (1) weight shuffling (topology preserved),
(2) per-neuron destination shuffling (out-degree preserved),
(3) full rewiring (configuration model, $n=500$ permutations
on $10^5$-edge subsample). Statistical significance:
$p = P(\alpha_{\mathrm{null}} \leq \alpha_{\mathrm{real}})$.

\subsection{Phase-space analysis}

Takens delay embedding~\cite{Takens1981} with $d=2,3$, $\tau=1$.
Correlation dimension $d_f$ from slope of $\ln C(r)$ vs $\ln r$.
Maximal Lyapunov exponent via Rosenstein algorithm~\cite{Rosenstein1993}.

%% ════════════════════════════════════════════════════════════
%%  III. SPECTRAL ANALYSIS
%% ════════════════════════════════════════════════════════════
\section{\label{sec:rmt}Spectral analysis and universal power law}

\subsection{Marchenko-Pastur analysis and synaptic attractors}

Figure~\ref{fig:marchenko} shows the eigenvalue spectrum of the
BBP L2/3 connectivity matrix. The bulk follows the MP distribution
($\lp = 89.0$), with $K = 20$ outlier eigenvalues
($\lmax = 720.9$, $\lmax/\lp = 8.1\times$) and level-spacing ratio
$\rGOE = 0.535 \approx 0.536$ (GOE prediction), confirming
Wigner-Dyson statistics.

For MICrONS ($\lp = 4.21$ after log-normalization): $K = 22$
outliers, $\lmax/\lp = 1{,}626\times$, $\rGOE = 0.654 > \mathrm{GOE}$.
The elevated $\rGOE$ in EM data may reflect directional asymmetry
(E/I balance) absent in the symmetrized simulation.

\subsection{Universal power law and AIC model selection}

Figure~\ref{fig:powerlaw} shows the ranked outlier eigenvalues
for both datasets and all four candidate models. Results:

\begin{equation}
  \lambda_k \sim k^\alpha, \quad \alpha = \begin{cases}
    -0.450 \pm 0.074 & \text{(BBP L2/3, }R^2=0.990\text{)} \\
    -0.422 \pm 0.062 & \text{(MICrONS V1, }R^2=0.984\text{)}
  \end{cases}
  \label{eq:alpha}
\end{equation}

Bootstrap 95\% CIs overlap ($[-0.539,-0.260]$ and $[-0.502,-0.267]$);
$H_0: \alpha_{\mathrm{BBP}} = \alpha_{\mathrm{MIC}}$ is not rejected.
AIC model selection (Fig.~\ref{fig:final}) confirms the power law
as uniquely preferred: $\Delta\mathrm{AIC} = 25.6$ (BBP) and
$23.0$ (MICrONS) vs.\ the next-best log-normal model. The harmonic
model ($1/k$, the $T{=}0$ quantum limit) is eliminated with
$\Delta\mathrm{AIC} > 219$.

%% ════════════════════════════════════════════════════════════
%%  IV. NULL MODEL VALIDATION
%% ════════════════════════════════════════════════════════════
\section{\label{sec:null}Null model validation}

\begin{table}[h]
\caption{\label{tab:null}Null model hierarchy (MICrONS V1).}
\begin{ruledtabular}
\begin{tabular}{lccc}
Model & $N_{\mathrm{out}}$ & $\lmax/\lp$ & $\alpha$ \\
\hline
Real (biological)         & 22 & 1626  & $-0.416$ \\
Null-1 (weights shuffled) & 22 & 1640  & $-0.416$ \\
Null-2 (destinations)     & 22 & 1620  & $-0.416$ \\
Null-3 (full rewiring)    & 20 &   74  & $-0.333$ \\
Null-3 distribution       & —  &    —  & $-0.141 \pm 0.009$ \\
\end{tabular}
\end{ruledtabular}
\end{table}

The number of outlier modes ($K\approx 20$) is topological
(preserved under Null-1,2), determined by the degree distribution.
The exponent $\alpha$ and energy ratio $\lmax/\lp$ are
biologically specific: both are destroyed by full rewiring
(Null-3), with $p < 0.001$ over 500 permutations
(Fig.~\ref{fig:bootstrap}).

%% ════════════════════════════════════════════════════════════
%%  V. LAMINAR GRADIENT
%% ════════════════════════════════════════════════════════════
\section{\label{sec:laminar}Laminar gradient and cross-dataset validation}

\begin{table}[h]
\caption{\label{tab:layers}Spectral properties by cortical layer.}
\begin{ruledtabular}
\begin{tabular}{lrrrrr}
Layer & $n$ & $K$ & $\alpha$ & $R^2$ & $\rGOE$ \\
\hline
BBP L2/3    &  8{,}173 & 20 & $-0.451$ & 0.992 & 0.558 \\
BBP L4      &  4{,}870 & 19 & $-0.500$ & 0.972 & 0.540 \\
BBP L5      &  5{,}824 & 19 & $-0.531$ & 0.973 & 0.515 \\
BBP L6      & 11{,}032 & 21 & $-0.531$ & 0.993 & 0.473 \\
\hline
MICrONS V1  & 50{,}948 & 22 & $-0.422$ & 0.984 & 0.654 \\
MICrONS L4  & 10{,}188 & 17 & $-0.557$ & 0.994 & —    \\
\end{tabular}
\end{ruledtabular}
\end{table}

$\alpha$ steepens monotonically with cortical depth
(Fig.~\ref{fig:laminar}), from $-0.451$ (L2/3) to $-0.531$
(L5/L6, $R^2 = 0.91$ for linear fit to depth), while $\rGOE$
decreases from $0.558$ toward Poisson. The layer-matched
comparison (BBP L4: $\alpha=-0.525$; MICrONS L4: $\alpha=-0.557$;
$\Delta\alpha = 0.032$, Fig.~\ref{fig:final}) confirms that the
difference between the full-dataset comparison ($\Delta\alpha=0.028$)
is primarily laminar, not species- or area-specific.

%% ════════════════════════════════════════════════════════════
%%  VI. THERMODYNAMIC DERIVATION
%% ════════════════════════════════════════════════════════════
\section{\label{sec:thermo}Thermodynamic derivation of EEG bands
         and dual UV cutoffs}

\subsection{Population integration time}

From patch-clamp recordings ($n=9$ cells, rat SSCx; 
Fig.~\ref{fig:tau}), median $\tauISI = 49.8\,\mathrm{ms}$.
The population integration time follows from the central limit
theorem applied to $k$ partially correlated spike trains:
\begin{equation}
  \taupop = \tauISI \cdot \sqrt{\kN} = 49.8\times\sqrt{0.15}
           = 19.3\,\mathrm{ms}.
  \label{eq:taupop}
\end{equation}
Three electrophysiological regimes are identified:
$\tauISI \in \{49.8, 8.5, 4.4\}\,\mathrm{ms}$ (rest, attention,
alert), corresponding to $\taupop \in \{19.3, 3.3, 1.7\}\,\mathrm{ms}$.

\subsection{Metabolic UV cutoff and EEG quantization}

The maximum number of non-overlapping ensembles before all $N$
neurons are simultaneously refractory:
\begin{equation}
  n_{\max} = \left\lfloor\frac{N}{k}\right\rfloor
           = \left\lfloor\frac{1}{\kN}\right\rfloor = 6.
  \label{eq:nmax}
\end{equation}
The $n$-th oscillatory mode has frequency:
\begin{equation}
  f_n(\mathrm{state}) = \frac{1}{n\cdot\tauISI(\mathrm{state})\cdot\sqrt{\kN}},
  \quad n = 1,\ldots,6.
  \label{eq:fn}
\end{equation}
The harmonic series $f_1:f_2:\cdots:f_6 = 1:1/2:\cdots:1/6$
is the invariant; $f_1$ scales with cognitive state through $\tauISI$.
At rest ($\taupop = 19.3\,\mathrm{ms}$) this yields:
$\{52, 26, 17, 13, 10, 9\}\,\mathrm{Hz}$, corresponding to the
canonical $\gamma$-through-$\theta$ bands.

\subsection{Dual UV cutoffs: Debye and Planck analogs}

Two distinct spectral limits emerge (Fig.~\ref{fig:planck}):

\textit{Ensemble cutoff} (Debye analog, state-dependent):
\begin{equation}
  \fUV^{\mathrm{ens}}(\mathrm{state}) = \frac{N/k-1}{\tauISI(\mathrm{state})}
  \label{eq:fUV_ens}
\end{equation}
At rest: $114\,\mathrm{Hz}$; under attention: $667\,\mathrm{Hz}$.

\textit{Neuronal absolute cutoff} (Planck analog, invariant):
\begin{equation}
  \fUV^{\mathrm{neural}} = \frac{1}{\tauabs\cdot\sqrt{\kN}}
  \approx 587\,\mathrm{Hz},
  \label{eq:fUV_neural}
\end{equation}
determined solely by the absolute refractory period
$\tauabs \approx 4.4\,\mathrm{ms}$ (measured). A remarkable
identity holds: at maximum alertness
($\tauISI \to \tauabs$), $f_1(\tauabs) = \fUV^{\mathrm{neural}}$
--- the fundamental mode coincides with the physical ceiling.

This invariant cutoff ($\approx 600\,\mathrm{Hz}$) is consistent
with empirical HFOs in human somatosensory cortex (407--926\,Hz
subdural~\cite{Hashimoto2009}) and thalamic VFOs
(1235--2632\,Hz~\cite{Hashimoto2009}), where the smaller thalamic
$\tauabs \approx 1\,\mathrm{ms}$ predicts
$\fUV^{\mathrm{thal}} = 587 \times (4.4/1) \approx 2580\,\mathrm{Hz}$,
in excellent agreement.

The dynamic EEG spectrum across states is shown in
Fig.~\ref{fig:spectrum}.

%% ════════════════════════════════════════════════════════════
%%  VII. PHASE SPACE
%% ════════════════════════════════════════════════════════════
\section{\label{sec:phase}Phase-space geometry of spectral attractors}

Takens delay embedding of the outlier eigenvalue sequence reveals
a quasi-1D attractor (Fig.~\ref{fig:takens}): correlation dimension
$d_f = 1.15$ (BBP, $R^2 = 0.98$) and $d_f = 1.07$ (MICrONS,
$R^2 = 0.99$). The slight excess above $d_f=1$ (a perfect power law)
reflects thermal fluctuations in $\kN$.

The maximal Lyapunov exponent is $\lambda_L = -0.224$ (BBP) and
$\lambda_L = -0.082$ (MICrONS), both negative (asymptotically
stable). The MICrONS value ($\lambda_L \approx -0.082$) lies
close to the critical point $\lambda_L = 0$, consistent with
the cortical criticality hypothesis~\cite{Beggs2003,Shew2013}
and suggesting that biological circuits operate closer to the
edge of chaos than deterministic models predict.

%% ════════════════════════════════════════════════════════════
%%  VIII. DISCUSSION
%% ════════════════════════════════════════════════════════════
\section{\label{sec:discussion}Discussion}

\subsection{A spectral invariant of cortical connectivity}

The exponent $\alpha \approx -0.43$ is conserved across two
independent datasets spanning different species (rat/mouse),
cortical areas (SSCx/V1), measurement modalities (model/EM),
and cortical layers (layer-matched $\Delta\alpha = 0.032$).
AIC model selection establishes it as a true power law, not
an exponential or harmonic decay. Its destruction by full rewiring
($p < 0.001$) establishes biological specificity beyond
degree distributions.

The laminar gradient ($\alpha$ steepening L2/3$\to$L5/L6,
$R^2=0.91$; $\rGOE$ decreasing toward Poisson) is consistent
with the functional specialization of cortical layers: superficial
layers perform distributed associative computation while deep
layers project topographically organized signals subcortically.

\subsection{Thermodynamic origin of EEG structure}

The derivation of $n_{\max}=6$ from metabolic refractoriness
provides a first-principles account of the six canonical EEG
bands without free parameters. The two UV cutoffs --- Debye
(state-dependent) and Planck (invariant, $\approx 587\,\mathrm{Hz}$)
--- predict the empirically observed ceiling on sustained cortical
oscillations. The thalamic prediction ($\approx 2580\,\mathrm{Hz}$)
from the same framework, matching observed VHFOs (1235--2632\,Hz),
constitutes an independent quantitative validation.

\subsection{Proximity to criticality}

The $\lambda_L = -0.082$ of MICrONS data (vs.\ $-0.224$ for
simulation) suggests real cortical circuits self-organize toward
the critical point~\cite{Beggs2003,Plenz2007}. The simulation's
stronger stability reflects intrinsic limitations of deterministic
models in capturing stochastic fluctuations driving criticality.

\subsection{Limitations}

The BBP dataset is a model with known idealizations.
Symmetrization $J_s = (J+J^T)/2$ discards directional information.
The MICrONS L4 extraction is approximate (coordinate-based, not
ground-truth layer assignment). The Takens embedding is limited
by small $K$ ($16$--$22$ outliers).

%% ════════════════════════════════════════════════════════════
%%  IX. CONCLUSIONS
%% ════════════════════════════════════════════════════════════
\section{\label{sec:conclusions}Conclusions}

The ranked eigenvalues of cortical synaptic weight matrices follow
a universal power law $\lambda_k \sim k^{-0.43}$ conserved between
simulation and electron microscopy data, across species, areas,
and layers, and destroyed by degree-preserving rewiring.
Formal AIC model selection ($\Delta\mathrm{AIC} > 23$) establishes
this as the preferred spectral model. The number of collective modes
($\sim 20$), their GOE-type level-spacing statistics, and the
quasi-1D phase-space geometry with $\lambda_L < 0$ (near-critical
in EM data) jointly characterize cortical connectivity as operating
in a stable but near-critical dynamical regime.
The metabolic constraint $\kN \approx 0.15$ governs all these
properties through the thermodynamic framework, including the
discrete EEG frequency bands, the Debye ensemble cutoff, and the
invariant Planck neuronal cutoff at $\approx 587\,\mathrm{Hz}$.

%% ════════════════════════════════════════════════════════════
%%  ACKNOWLEDGMENTS
%% ════════════════════════════════════════════════════════════
\begin{acknowledgments}
The author thanks the Blue Brain Project and the Open Brain
Institute for publicly releasing the neocortical connectivity
data via AWS Open Data, and the MICrONS Consortium for the
electron microscopy reconstruction. All analyses used open-source
software (Python, NumPy, SciPy, h5py, Matplotlib) and
publicly available datasets.
\end{acknowledgments}

\section*{Data Availability}
All datasets are publicly available:
BBP connectivity matrices at \texttt{s3://openbluebrain/}
(no credentials required); MICrONS V1 at the same repository.
Analysis code available at \texttt{[GitHub URL]}.

%% ════════════════════════════════════════════════════════════
%%  FIGURES
%% ════════════════════════════════════════════════════════════

\begin{figure}[t]
  \includegraphics[width=\columnwidth]{marchenko_pastur_completo}
  \caption{\label{fig:marchenko}
    \textbf{Marchenko-Pastur analysis — BBP L2/3.}
    ($n = 8{,}173$ neurons, $974{,}129$ synapses.)
    \textit{Left}: Top 30 eigenvalues; orange bars exceed the MP
    upper edge $\lp = 89.0$ (dashed red), identifying 20 synaptic
    attractors. \textit{Right}: Same data on log scale confirming
    $\lmax/\lp = 8.1\times$.}
\end{figure}

\begin{figure}[t]
  \includegraphics[width=\columnwidth]{BBP_vs_MICrONS_final}
  \caption{\label{fig:powerlaw}
    \textbf{Universal power-law spectrum of synaptic attractors.}
    BBP L2/3 (red, simulation, rat SSCx) vs.\ MICrONS V1
    (green, electron microscopy, mouse V1).
    \textit{Upper row}: Ranked eigenvalue spectra with power-law
    fits ($\alpha = -0.450$, $R^2 = 0.990$ and $\alpha = -0.416$,
    $R^2 = 0.986$). Both datasets are well above their respective
    MP upper edges (dotted lines).
    \textit{Lower row}: Eigenvalue spacings showing level-spacing
    ratios $\rGOE = 0.535 \approx \mathrm{GOE}$ (BBP) and
    $\rGOE = 0.654 > \mathrm{GOE}$ (MICrONS).}
\end{figure}

\begin{figure}[t]
  \includegraphics[width=\columnwidth]{bootstrap_alpha}
  \caption{\label{fig:bootstrap}
    \textbf{Bootstrap analysis and null model validation.}
    \textit{Upper left/center}: Bootstrap distributions of $\alpha$
    ($n = 1{,}000$) for BBP and MICrONS with 95\% CIs.
    CIs overlap, confirming $H_0: \alpha_{\mathrm{BBP}} =
    \alpha_{\mathrm{MIC}}$.
    \textit{Upper right}: Null model distribution ($n = 500$
    full-rewiring permutations, $\bar\alpha = -0.141$),
    completely separated from real values ($p < 0.001$).
    \textit{Lower left}: Laminar gradient with bootstrap CIs.
    \textit{Lower center/right}: Summary statistics.}
\end{figure}

\begin{figure}[t]
  \includegraphics[width=\columnwidth]{layers_spectral}
  \caption{\label{fig:laminar}
    \textbf{Laminar gradient of the spectral exponent — BBP HEX0.}
    \textit{Left}: $\alpha$ by layer; dashed line = MICrONS reference.
    \textit{Center}: Level-spacing ratio $\rGOE$ by layer, with
    GOE (dashed) and Poisson (dotted) references.
    \textit{Right}: $\alpha$ vs.\ cortical depth, linear trend
    $R^2 = 0.91$. Superficial layers (L2/3) have flatter spectra
    ($\alpha = -0.451$) and stronger inter-attractor repulsion
    ($\rGOE > \mathrm{GOE}$); deep layers (L5/L6) have steeper
    spectra ($\alpha = -0.531$) and more independent modes
    ($\rGOE < \mathrm{GOE}$).}
\end{figure}

\begin{figure}[t]
  \includegraphics[width=\columnwidth]{planck_cognitivo}
  \caption{\label{fig:planck}
    \textbf{Cognitive Planck quantization and UV cutoffs.}
    \textit{Upper left}: Bose-Einstein occupancy vs.\ frequency,
    comparing classical Rayleigh-Jeans divergence with the cognitive
    Planck distribution ($\fUV = 114\,\mathrm{Hz}$ at rest).
    Vertical lines mark the six quantized EEG modes.
    \textit{Upper right}: Ensemble UV cutoff $\fUV^{\mathrm{ens}}$
    vs.\ activation fraction $\kN$.
    \textit{Lower left}: Number of sustainable activation cycles
    before metabolic collapse; annotation at $n_{\max} = N/k = 6$.
    \textit{Lower right}: Quantized modes with Planck occupancy
    (bars) and frequencies (green dots).}
\end{figure}

\begin{figure}[t]
  \includegraphics[width=\columnwidth]{UV_catastrophe_v2}
  \caption{\label{fig:uv}
    \textbf{Cognitive UV catastrophe and phase diagram.}
    \textit{Upper left}: Population integration time
    $\taupop = \tauISI\sqrt{\kN}$ vs.\ activation; $\taupop = 19.3\,\mathrm{ms}$
    at $\kN = 0.15$.
    \textit{Upper right}: Dominant frequency vs.\ activation,
    with labeled cognitive states (sustained attention:
    $\kN \approx 0.30$; epileptic crisis: $\kN \to 1$).
    \textit{Lower left}: Cognitive efficiency
    (capacity/metabolic cost); $\kN = 0.15$ lies in the high-efficiency
    plateau.
    \textit{Lower right}: Cognitive phase diagram; shaded zone
    marks the UV catastrophe region where population computation
    collapses.}
\end{figure}

\begin{figure}[t]
  \includegraphics[width=\columnwidth]{takens_phase_space}
  \caption{\label{fig:takens}
    \textbf{Phase-space structure via Takens delay embedding.}
    Columns: BBP L2/3 (left) and MICrONS V1 (right).
    \textit{Top}: Raw eigenvalue spectra.
    \textit{Second row}: 2D phase-space trajectories
    $(\lambda_i, \lambda_{i+1})$, colored by rank.
    Both datasets trace smooth, quasi-1D curves characteristic
    of low-dimensional attractors.
    \textit{Third row}: 3D attractors confirming planar geometry.
    \textit{Bottom}: Correlation dimension ($d_f = 1.15$ and $1.07$;
    $R^2 > 0.98$). Maximal Lyapunov exponents:
    $\lambda_L = -0.224$ (BBP, stable) and $-0.082$ (MICrONS,
    near-critical).}
\end{figure}

\begin{figure}[t]
  \includegraphics[width=\columnwidth]{dynamic_spectrum}
  \caption{\label{fig:spectrum}
    \textbf{Dynamic EEG spectrum across cognitive states.}
    \textit{Left}: Frequencies $f_n = 1/(n\taupop)$ for $n=1,\ldots,6$
    under three regimes (rest: blue; attention: orange; alert: red).
    $n_{\max} = 6$ is invariant; only the frequency scale changes.
    Shaded bands: clinical EEG ranges.
    \textit{Right}: $f_1$ (gamma) and $f_6$ (theta) vs.\ $\tauISI$
    on log-log axes. Vertical lines mark the three measured
    refractory regimes. The invariant neuronal cutoff
    $\fUV^{\mathrm{neural}} \approx 587\,\mathrm{Hz}$
    (where $f_1(\tauabs) = \fUV^{\mathrm{neural}}$)
    is off-scale to the right.}
\end{figure}

\begin{figure}[t]
  \includegraphics[width=\columnwidth]{final_tests}
  \caption{\label{fig:final}
    \textbf{AIC model selection and layer-matched validation.}
    \textit{Left}: $\Delta\mathrm{AIC}$ for four candidate models
    (BBP: red; MICrONS: green). The power law wins in both datasets;
    $\Delta\mathrm{AIC} > 10$ (dashed line) denotes strong evidence.
    \textit{Center}: Visual fits for BBP L23 confirming that the
    power law (solid) captures the curvature missed by the Boltzmann
    (dashed) and harmonic (dotted) models.
    \textit{Right}: Layer-matched comparison. BBP L4: $\alpha=-0.525$;
    MICrONS L4: $\alpha=-0.557$; $\Delta\alpha = 0.032$, within
    bootstrap uncertainty. MICrONS total ($-0.422$, dashed) aligns
    with BBP L2/3, consistent with the predominantly superficial
    sampling of the MICrONS volume.}
\end{figure}

%% ── Supplemental figures ─────────────────────────────────────

\begin{figure}[t]
  \includegraphics[width=\columnwidth]{tau_ISI_refinado}
  \caption{\label{fig:tau}
    \textbf{Population integration time — electrophysiological
    validation.}
    \textit{Left}: Distribution of ISI minimum (blue) and ISI p5
    (orange) per cell ($n=9$, rat SSCx, IDRest protocol,
    highest-current sweep). Median ISI$_{\min} = 8.5\,\mathrm{ms}$
    (attention regime); median ISI$_{\mathrm{p5}} = 29.2\,\mathrm{ms}$.
    Red dotted line: TCSV $\tau = 20\,\mathrm{ms}$, derived from
    median ISI $= 49.8\,\mathrm{ms}$ via
    $\taupop = \tauISI\sqrt{\kN} = 19.3\,\mathrm{ms}$.
    \textit{Right}: Predicted EEG band frequencies for three
    electrophysiological regimes vs.\ clinical references.}
\end{figure}

\begin{figure}[t]
  \includegraphics[width=\columnwidth]{h_neural_volume}
  \caption{\label{fig:hneural}
    \textbf{Synaptic volume and informational density.}
    MICrONS EM ($n = 11{,}554{,}213$ synapses). 
    \textit{Upper left}: Log-normal distribution of synaptic volumes
    ($\mu = 8.17$, $\sigma = 1.03$; Shannon entropy $H = 2.02\,\mathrm{bits}$).
    \textit{Lower panels}: $h_{\mathrm{neural}}$ derived via
    $E \cdot \tau / H$ as a function of synaptic volume.
    The TCSV postulate $\hneu = 10^{-13}\,\mathrm{J\cdot s}$
    corresponds to the activation entropy
    $H_{\mathrm{act}} = \ln\binom{100}{15} = 58\,\mathrm{bits}$
    of the microcircuit ensemble, not to the volume-encoded
    structural information. The log-normal CV $\approx 1.04$
    confirms the hyperbolic geometry of the semantic manifold.}
\end{figure}

%% ════════════════════════════════════════════════════════════
%%  BIBLIOGRAPHY
%% ════════════════════════════════════════════════════════════
\begin{thebibliography}{99}

\bibitem{Mountcastle1997}
V.~B. Mountcastle,
Brain \textbf{120}, 701 (1997).

\bibitem{Attwell2001}
D.~Attwell and S.~B. Laughlin,
J.\ Cereb.\ Blood Flow Metab.\ \textbf{21}, 1133 (2001).

\bibitem{Mehta2004}
M.~L. Mehta,
\emph{Random Matrices}, 3rd ed.\ (Elsevier, 2004).

\bibitem{Mastrogiuseppe2018}
F.~Mastrogiuseppe and S.~Ostojic,
Neuron \textbf{99}, 609 (2018).

\bibitem{Markram2015}
H.~Markram \emph{et al.},
Cell \textbf{163}, 456 (2015).

\bibitem{MICrONS2021}
MICrONS Consortium \emph{et al.},
bioRxiv 2021.07.28.454025 (2021).

\bibitem{Bhatt2009}
D.~L. Bhatt \emph{et al.},
J.\ Neurosci.\ \textbf{29}, 2065 (2009).

\bibitem{Atas2013}
Y.~Y. Atas \emph{et al.},
Phys.\ Rev.\ Lett.\ \textbf{110}, 084101 (2013).

\bibitem{Takens1981}
F.~Takens, in \emph{Dynamical Systems and Turbulence},
Lecture Notes in Mathematics \textbf{898} (Springer, 1981), p.~366.

\bibitem{Rosenstein1993}
M.~T. Rosenstein, J.~J. Collins, and C.~J. De~Luca,
Physica D \textbf{65}, 117 (1993).

\bibitem{Beggs2003}
J.~M. Beggs and D.~Plenz,
J.\ Neurosci.\ \textbf{23}, 11167 (2003).

\bibitem{Plenz2007}
D.~Plenz and T.~C. Thiagarajan,
Trends Neurosci.\ \textbf{30}, 101 (2007).

\bibitem{Shew2013}
W.~L. Shew and D.~Plenz,
Neuroscientist \textbf{19}, 88 (2013).

\bibitem{Hashimoto2009}
I.~Hashimoto \emph{et al.},
Epilepsia \textbf{50}, 2049 (2009).

\bibitem{Pribylova2024}
L.~P\v{r}ibylov\'a \emph{et al.},
Network Neurosci.\ \textbf{8}, 293 (2024).

\end{thebibliography}

\end{document}
